\section{Use Case and Observations} \label{sec:usecase}
We applied ODDToolkit within the ontology engineering lifecycle of a Flemish government project for environmental reporting. In earlier projects, ontologies used for publication were often retrofitted to data models that had already been shaped by implementation decisions. For the IED~(Industrial Emissions Directive) and the IEPR regulation~\cite{EU2024_1244}, we instead started from a new ontology intended to model industrial processes and their emissions, with implementation artefacts generated around that conceptual model. This change in approach required a new workflow for ontology engineers and developers to understand each other's artefacts.

The ontology reused SOSA~\cite{janowicz2019sosa}, P-PLAN~\cite{garijo2012pplan}, PROV-O, and GeoSPARQL, which provided an interoperable modelling basis but also introduced abstractions and predicate names that were not immediately aligned with the development team's implementation vocabulary, a challenge consistent with prior observations about linked data adoption~\cite{9963753,avila2020applications}. This was particularly visible for developers who needed to understand the IED process model from technology-facing artefacts rather than from ontology axioms directly. Furthermore, developers were not familiar with the openness of linked data, resulting in a preference for more traditional artefacts.

The ontology, SKOS concept scheme, and YAML configuration were maintained in version control and executed in continuous integration so artefacts were regenerated and communicated on each ontology change. The SKOS concept scheme was jointly authored by ontology engineers and lead developers, capturing naming decisions that could not be inferred from the ontology alone; for instance, \texttt{sosa:isHostedBy} was mapped to \texttt{location} because the only relevant hosting relationship for an emission point in this domain is between the emission point and its operation location. Design decisions such as the use of surrogate keys for versioning were captured in the YAML configuration.

In practice, the generated Java and TypeScript models, as well as the SQL DDL were used as implementation starting points. The diagrams supported project documentation, and the SHACL shapes fed validation and synthetic-data workflows. While the project is ongoing, the generated artefacts serve as technology-familiar entry points to the IEPR data model, complementing the formal ontology definitions rather than requiring developers to start from ontology axioms alone. This observation is not based on a comparative evaluation or developer interviews, but it illustrates how ODDToolkit can keep implementation, documentation, and validation artefacts aligned around a shared ontology.
